% Deterministically generated from the frozen audited Markdown source.
% Responsible author: Carptopus.
\documentclass[11pt]{article}
\usepackage[a4paper,margin=25mm]{geometry}
\usepackage[T1]{fontenc}
\usepackage{lmodern,microtype}
\usepackage{amsmath,amssymb,amsthm,mathtools}
\renewcommand{\sfdefault}{\rmdefault}
% The frozen source uses math sans only for three local polynomial labels.
% Render those labels upright because the local offline cache lacks lmss10.pfb.
\renewcommand{\mathsf}[1]{\mathrm{#1}}
\usepackage{needspace}
\usepackage{xcolor,xurl,fancyhdr}
\usepackage[colorlinks=true,linkcolor=blue!55!black,citecolor=blue!55!black,urlcolor=blue!65!black]{hyperref}

\pagestyle{fancy}
\fancyhf{}
\fancyhead[L]{Carptopus}
\fancyhead[R]{Three-branch trees and the CSF}
\fancyfoot[C]{\thepage}
\setlength{\headheight}{14pt}
\setlength{\parindent}{0pt}
\setlength{\parskip}{0.42em}
\setlength{\emergencystretch}{4em}
\allowdisplaybreaks
\urlstyle{same}

\title{Trees with three branch vertices are determined\\
by their chromatic symmetric functions}
\author{Carptopus\\\href{mailto:carptopus@163.com}{carptopus@163.com}}
\date{17 September 2026}

\hypersetup{
  pdftitle={Trees with three branch vertices are determined by their chromatic symmetric functions},
  pdfauthor={Carptopus},
  pdfsubject={Chromatic symmetric reconstruction of trees with exactly three branch vertices},
  pdfkeywords={chromatic symmetric function, tree reconstruction, branch vertices, spider, subtree polynomial, path sequence}
}

\begin{document}
\maketitle

Internal draft, 17 September 2026. Internal pre-writing and full-manuscript mathematical reviews
completed; external peer review pending.
\begin{abstract}

Stanley asked whether the chromatic symmetric function distinguishes nonisomorphic trees. The
conjecture is known for several structured families, including spiders and, more recently, all trees
with exactly two vertices of degree at least three. We prove the next complete branch-vertex layer:
every finite tree with exactly three vertices of degree at least three is determined, among all
finite simple graphs, by its chromatic symmetric function.

The three branch vertices lie on a path. We combine a power-sum derivative of an alternating
trunk deletion sum with the path sequence and the three-leaf part of the subtree polynomial. The
derivative first recovers the smallest terminal starlike block and then produces a sequence of
endpoint events. Equal terminal orders, rooted paths, simultaneous events, and the two possible
scalar-layout exchanges require separate arguments. Their residual ambiguities reduce to explicit
polynomial factorizations or threshold-moment systems, all of which are nonsingular or have
classifiable degenerate cases. The resulting reconstruction is uniform in the trunk lengths, twig
lengths, and branch degrees. Exact computations are used only as regression and boundary checks,
not as substitutes for the general proof.

\textbf{Keywords:} chromatic symmetric function; tree reconstruction; branch vertices; spider; subtree
polynomial; path sequence

\end{abstract}

\section{Introduction}


For a finite simple graph $G$, Stanley's chromatic symmetric function is

\begin{equation*}
X_G=\sum_{\kappa}\prod_{v\in V(G)}x_{\kappa(v)},
\end{equation*}


where the sum ranges over all proper colorings of $G$ by positive integers. Stanley asked whether
$X_T=X_{T'}$ for two trees forces $T\cong T'$ [1]. The question remains open in general.

Several infinite classes are known. Martin, Morin and Wagner proved that spiders are determined and
extracted path, degree and subtree data from $X_T$ [2]. Crew showed that $X_T$ determines the order
of the trunk and the multiset of twig lengths [3]. Loebl and Sereni treated caterpillars [4], Huryn
and Chmutov treated 2-spiders [5], and later work established proper $q$-caterpillars and all trees
of diameter at most five [6,7]. Recent multiplicity-isolation criteria cover further proper trees
[10,11]. These families intersect the branch-vertex filtration but do not contain a whole new
layer of it.

Wang, Yu and Zhang proved that every tree with exactly two vertices of degree at least three is
determined by its chromatic symmetric function [8]. Their proof combines a weighted
deletion--contraction construction along the path joining the two branch vertices with a
reconstruction lemma for pairs of spiders. They remarked that the method might extend to three
branch vertices. We carry out that extension, including the path-rooting and simultaneous-event
phenomena that do not occur in the two-branch layer.

Call a vertex of degree at least three a \textbf{branch vertex}. Our main result is the following.

\Needspace{0.13\textheight}
\subsection*{Theorem 1.1}


Let $T$ be a finite simple tree with exactly three branch vertices. If a finite simple graph $F$
satisfies $X_F=X_T$, then $F\cong T$.

Together with the known zero-, one-, and two-branch cases, this gives the following attributed
corollary.

\Needspace{0.13\textheight}
\subsection*{Corollary 1.2}


Every finite tree with at most three branch vertices is determined, among finite simple graphs, by
its chromatic symmetric function.

The new content of the paper is the exact-three layer. Four branch vertices introduce branched
cores and new placement types, so Theorem 1.1 does not imply the next layer and does not resolve
Stanley's conjecture.

The proof has three ingredients. First, an alternating deletion sum over the trunk gives a family
of power-sum derivatives whose lowest nonzero member recovers the smallest terminal block. Second,
the path sequence fixes the two trunk gaps once the three rooted blocks are known. Third, the path
sequence and the three-leaf subtree polynomial eliminate the allocation exchanges that remain when
two derivative events occur at the same order. The technically delicate case is a smallest terminal
block that is an unrooted path: the chromatic symmetric function of the block does not remember the
position of its attachment vertex. Sections 5--7 give a uniform rooting argument and classify the
only two possible changes of scalar layout.

\section{Preliminaries and notation}


We use the power sums $p_1,p_2,\ldots$. Stanley's expansion is

\begin{equation*}
X_G=\sum_{S\subseteq E(G)}(-1)^{|S|}p_{\lambda(S)},
\tag{2.1}
\end{equation*}


where $\lambda(S)$ is the partition of $|V(G)|$ formed by the component orders of $(V(G),S)$
[1]. Consequently $X_G$ determines the vertex and edge counts, the girth, the degree sequence of a
tree, and the number of one-edge cuts of every component-order type [1,2,9]. We also use that for
trees the chromatic symmetric function determines the path sequence and the subtree polynomial;
the precise equivalence with generalized degree and subtree data is recorded in [9].

For a tree $T$, let $P$ be the smallest subtree containing all branch vertices. Crew calls $P$ the
\textbf{trunk} [3]. A \textbf{twig} is a component path left after removing the trunk edges, with its trunk
endpoint retained as root; its length is its number of edges. Crew proved that $X_T$ determines
$|P|$ and the multiset of all twig lengths [3].

When $T$ has exactly three branch vertices, the trunk is a path. Orient it temporarily from left to
right and write its branch vertices as $u,v,w$. Delete the two trunk segments while retaining their
endpoints. The rooted starlike blocks at $u,v,w$ are denoted by $A,B,C$, and

\begin{equation*}
a=|A|,\qquad b=|B|,\qquad c=|C|.
\end{equation*}


Let $\alpha,\beta\geq1$ be the distances from $u$ to $v$ and from $v$ to $w$. Thus

\begin{equation*}
|T|=a+b+c+\alpha+\beta-2.
\tag{2.2}
\end{equation*}


The endpoint blocks $A,C$ have at least two twigs; $B$ has at least one. Reversing the entire
display is immaterial.

For a rooted starlike block $D$ with arm-length multiset $\Lambda_D$, set

\begin{equation*}
f_D(x)=\sum_{z\in V(D)}x^{\mathrm{dist}(z,\mathrm{root}(D))}
      =1+\sum_{\ell\in\Lambda_D}h_\ell(x),
\qquad h_\ell(x)=x+\cdots+x^\ell.
\tag{2.3}
\end{equation*}


The arm-sum map $\Lambda\mapsto\sum_{\ell\in\Lambda}h_\ell$ is injective: the coefficient of
$x^i$ is the number of arms of length at least $i$.

We write

\begin{equation*}
W_T(x)=\sum_{\{r,s\}\subseteq V(T)}x^{\mathrm{dist}(r,s)}
\tag{2.4}
\end{equation*}


for the path sequence, with unordered pairs of distinct vertices. Let $S_T(x,y)$ be the subtree
polynomial in which a connected subtree with $e$ edges and $j$ leaves contributes $x^e y^j$.
The coefficient

\begin{equation*}
J_T(x)=[y^3]S_T(x,y)
\tag{2.5}
\end{equation*}


is the three-leaf part. Both (2.4) and (2.5) are determined by $X_T$ [2,9]. For an arm multiset
$U$, write $e_i(U)$ for the $i$th elementary symmetric polynomial in
$\{h_\ell(x):\ell\in U\}$. A spider with arms $U$ has

\begin{equation*}
J_U(x)=e_3(U).
\tag{2.6}
\end{equation*}


We use two established reconstruction results. First, a spider is determined by its chromatic
symmetric function [2]. Second, a tree with exactly two branch vertices is determined by its
chromatic symmetric function [8]. Whenever a rooted starlike block has at least three arms, its
root is the unique branch vertex, so unrooted reconstruction also recovers the root. A two-arm
block is an unrooted path, and its root position must be recovered separately.

\section{The trunk complement transform}


Let the vertices of the full trunk be $p_0,p_1,\ldots,p_m$, with the three branch vertices at
$p_0,p_h,p_m$. Define the alternating trunk sum

\begin{equation*}
\mathcal A_P(T)=
\sum_{\varnothing\neq I\subseteq E(P)}(-1)^{|I|-1}X_{T-I}.
\tag{3.1}
\end{equation*}


The weighted deletion--contraction formula of Wang, Yu and Zhang, together with the recoverable
trunk order and twig sequence, shows that $\mathcal A_P(T)$ is determined by $X_T$ [8].

Cutting the terminal trunk edge at the left or right leaves a terminal starlike block; call these
$L$ and $R$. Put

\begin{equation*}
u=\min\{|L|,|R|\},\qquad n=|T|,\qquad k=n-u.
\end{equation*}


\Needspace{0.13\textheight}
\subsection*{Lemma 3.1 (minimal terminal extraction)}


The value $u$ and the sum of the chromatic symmetric functions of all terminal blocks of order $u$
are determined by $X_T$. More precisely,

\begin{equation*}
(-1)^{k-1}\frac{\partial\mathcal A_P(T)}{\partial p_k}
=\sum_{E\in\{L,R\},\ |E|=u}X_E.
\tag{3.2}
\end{equation*}


\paragraph{Proof}


Fix a connected induced $k$-vertex subtree $H$. In the derivative of (3.1), group all terms that
select $H$ as the $p_k$ component. Let $B_H$ be the trunk edges with one endpoint in $H$, and
$C_H$ the trunk edges with neither endpoint in $H$. After removing $V(H)$ the integrand is
independent of the selected edges in $B_H$. Inclusion--exclusion over $B_H$ therefore gives one if
$B_H\neq\varnothing$ and zero if $H$ contains the entire trunk.

Every twig is shorter than $k$: after deleting a twig, a terminal block of order at least $u$ and
another branch vertex remain. Thus $H$ meets the trunk. If it does not contain the whole trunk,
$T-V(H)$ contains a complete terminal block and has order $u$. It must therefore equal precisely a
terminal block of minimum order. This proves (3.2). For $s<u$ the same derivative with
$p_{n-s}$ vanishes, whereas at $s=u$ the right side of (3.2) is nonzero. Hence $u$ is recovered as
the first nonzero index. $\square$

To continue past the first coefficient, define

\begin{equation*}
q_\ell(z)=\sum_{r=0}^{\ell}X_{P_r}z^r,\qquad X_{P_0}=1,
\quad
Q_v(z)=\prod_{\ell\in\Lambda_v}q_\ell(z).
\tag{3.3}
\end{equation*}


For a proper nonempty trunk interval $[i,j]$, let $L_i$ and $R_j$ be the full left and right
complement blocks; an empty side has chromatic symmetric function one. Let $\lambda$ be the maximum
twig length and set

\begin{equation*}
E_s(T)=(-1)^{n-s-1}\frac{\partial\mathcal A_P(T)}{\partial p_{n-s}}.
\end{equation*}


\Needspace{0.13\textheight}
\subsection*{Lemma 3.2 (interval complement identity)}


Truncated in degrees $s<n-\lambda$,

\begin{equation*}
\sum_sE_s(T)z^s=
\sum_{[i,j]}
X_{L_i}X_{R_j}z^{|L_i|+|R_j|}
\prod_{v\in\{p_0,p_h,p_m\}\cap[i,j]}Q_v(z),
\tag{3.4}
\end{equation*}


where the sum is over proper nonempty trunk intervals.

\paragraph{Proof}


A connected $(n-s)$-vertex subtree cannot lie wholly inside a twig when $n-s>\lambda$, so its
intersection with the trunk is a nonempty interval. The full-trunk interval cancels by the same
inclusion--exclusion as in Lemma 3.1. For a proper interval, the outside trunk pieces force the
factors $X_{L_i}X_{R_j}$. At each branch vertex inside the interval, the connected subtree retains
an arbitrary initial segment of every twig, leaving a terminal path forest enumerated by $Q_v$.
The choices are independent, giving (3.4). $\square$

Suppose henceforth that the left endpoint block is the unique smaller endpoint, $a<c$, and that its
rooted arm multiset is known. Let $A^t$ denote $A$ with an extra arm of length $t$ at its root,
where $A^0=A$. Dividing (3.4) by the known pruning series shows that the initial coefficients are the
hypothetical chain

\begin{equation*}
X_{A^0},X_{A^1},\ldots,X_{A^{\alpha-1}}.
\end{equation*}


The first departure has order

\begin{equation*}
s_0=\min(a+\alpha,c)
\tag{3.5}
\end{equation*}


and equals

\begin{equation*}
\Delta_{s_0}=
\begin{cases}
-X_{A^\alpha},&a+\alpha<c,\\
 X_C,&c<a+\alpha,\\
 X_C-X_{A^\alpha},&c=a+\alpha.
\end{cases}
\tag{3.6}
\end{equation*}


All three alternatives lie in the safe range of (3.4). Indeed, using (2.2), $n-s_0$ is larger
than every twig length. If the last line cancels because $C\cong A^\alpha$, the next coefficient is

\begin{equation*}
\Delta_{c+1}=-X_{A^{\alpha+1}}-p_1X_C
+\mathbf1_{\beta\geq2}X_{C^1}
+\mathbf1_{b=2}X_{A\mathbin{-}B}.
\tag{3.7}
\end{equation*}


Its principal specialization divided by $q(q-1)^{c-1}$ is one of $-q,-1,1-2q$, whereas an
ordinary left or right first event gives $-(q-1)$ or $q-1$ and an equal-order difference gives zero.
Thus seamless cancellation cannot hide the event across two candidates.

Once $\alpha$ is known, put

\begin{equation*}
D=A\mathbin{-}_{\alpha}B,\qquad d=|D|=a+b+\alpha-1.
\tag{3.8}
\end{equation*}


Subtracting the true finite left chain from (3.4), the next event at $\min(c,d)$ is respectively
$X_C$, $X_D$, or $X_C+X_D$. The remainder of the proof separates the equal-order sum and restores
root positions.

\section{Recovering the trunk split}


Once the rooted blocks $A,B,C$ are known, no further high-order cut data are needed.

\Needspace{0.13\textheight}
\subsection*{Lemma 4.1 (trunk split rigidity)}


Fix the rooted blocks $A,B,C$ and the total trunk length $m=\alpha+\beta$. The path sequence
determines $\alpha,\beta$ up to simultaneous left--right reversal.

\paragraph{Proof}


Let $g_B=f_B-1$. The part depending on the position $\alpha$ of the middle root is

\begin{equation*}
W_{T_\alpha}(x)=W_0(x)+g_B(x)H_\alpha(x),
\tag{4.1}
\end{equation*}


where $W_0$ is independent of $\alpha$ and

\begin{equation*}
H_\alpha=x^\alpha f_A+x^{m-\alpha}f_C
+\sum_{i=1}^{\alpha-1}x^i+\sum_{i=1}^{m-\alpha-1}x^i.
\tag{4.2}
\end{equation*}


Since $B$ has a twig, $g_B\neq0$. If $\alpha<\gamma$ and $H_\alpha=H_\gamma$, put
$r=\gamma-\alpha$ and $\delta=m-\gamma$. Direct subtraction gives

\begin{equation*}
0=(1+x+\cdots+x^{r-1})
\bigl[-x^\alpha G_A+x^\delta G_C\bigr],
\tag{4.3}
\end{equation*}


where

\begin{equation*}
G_D=1-(1-x)f_D=(1-|\Lambda_D|)x+
\sum_{\ell\in\Lambda_D}x^{\ell+1}.
\end{equation*}


The least nonzero degree of each endpoint $G_D$ is one. Equation (4.3) forces
$\alpha=\delta=m-\gamma$ and then $G_A=G_C$. The latter equality recovers equal rooted arm
multisets, and the two placements differ only by reversing the tree. $\square$

\section{Equal terminal orders}


Assume $a=c$. Lemma 3.1 gives $X_A+X_C$. Let $\Pi_0$ retain only power-sum monomials not divisible
by $p_1$. At complement order $a+1$, every pruning correction contains an isolated vertex, while an
endpoint with gap at least two can be extended by one trunk vertex. Hence

\begin{equation*}
\Pi_0E_{a+1}=
\Pi_0\bigl(
\mathbf1_{\alpha\geq2}X_{A^1}+
\mathbf1_{\beta\geq2}X_{C^1}
\bigr).
\tag{5.1}
\end{equation*}


For an $M$-vertex spider with threshold numbers
$t_i=\#\{\ell\in\Lambda:\ell\geq i\}$, the one-edge-cut counts satisfy

\begin{equation*}
d_i=t_i+t_{M-i}\quad(i<M/2),
\qquad d_{M/2}=t_{M/2}\quad(M\text{ even}).
\tag{5.2}
\end{equation*}


These counts determine every spider with at least three arms. If a long arm exists it is unique;
two putative long-arm solutions to (5.2) could differ only by replacing
$\{k-1,M-k\}$ with $\{\ell-1,M-\ell\}$, leaving no room for a third positive arm. Thus the only
fold ambiguity is a two-arm path. Since $A^1$ and $C^1$ have at least three arms and contain a
distinguished new arm of length one, (5.1) recovers every active endpoint.

If both gaps are at least two, the two extended endpoint blocks give, for $i\geq2$,

\begin{equation*}
d_i=T_i+T_{a-i},\qquad e_i=T_i+T_{a+1-i},
\tag{5.3}
\end{equation*}


where $T_i=t_i(A)+t_i(C)$. Successive differences recover all $T_i$. The merged arm multiset and
$X_A+X_C$ then recover the unordered rooted pair $\{A,C\}$ as follows.

\Needspace{0.13\textheight}
\subsection*{Lemma 5.1 (equal-order endpoint pair)}


Let $A,C$ be rooted starlike blocks of equal order. Their chromatic-symmetric-function sum, their
forest path sequence, degree sequence, and merged arm multiset determine the unordered rooted pair.

\paragraph{Proof}


If both blocks have at least three arms, apply the two-spider reconstruction lemma of [8]. If exactly
one is a path, its unrooted type is the unique path of the common order; subtract its chromatic
symmetric function and recover the other spider, then recover the path root by subtracting the
spider arms from the merged multiset. If both are paths, each rooted path consists of two positive
arms whose lengths sum to $a-1$. The involution $r\mapsto a-1-r$ uniquely pairs the four-arm
multiset, including the fixed-point case. $\square$

If exactly one gap is at least two, (5.1) first recovers the endpoint on the long-gap side. The
other endpoint is then recovered unless it is a path. In that remaining case, let $a(x),b(x),c(x)$
be the arm sums at the adjacent path endpoint, middle block, and already known endpoint; let
$u=a+b$ and let the other gap be $\beta\geq2$. The allocation-dependent part of the path sequence is

\begin{equation*}
V(a,b)=e_2(A)+e_2(B)+x(1+a)(1+b)
+x^\beta(1+b)f_C+x^{\beta+1}(1+a)f_C.
\tag{5.4}
\end{equation*}


For another allocation $a',b'$ with $a+b=a'+b'=u$,

\begin{equation*}
V(a,b)-V(a',b')=(a-a')(1-x)
\bigl[a+a'-u-x^\beta f_C\bigr].
\tag{5.5}
\end{equation*}


The last factor evaluates at $x=1$ to $-|B|\neq0$. Hence $a=a'$, and injectivity of the arm-sum
encoding fixes the rooted path.

It remains to treat $\alpha=\beta=1$. Let $E(z)=\sum_sE_s z^s$ and $Q=Q_AQ_BQ_C$ in
(3.3). Before the middle block can be crossed, formal division of (3.4) gives

\begin{equation*}
z^{-a}E(z)/Q(z)=X_A/Q_A(z)+X_C/Q_C(z)
\pmod {z^{\min(a,b)}}.
\tag{5.6}
\end{equation*}


In the coefficient of $z^k$, the only term that can contribute the power-sum monomial $p_ap_k$
uses the unique one-component term of $X_A$ or $X_C$ and the unique one-component term of the
inverse pruning series. Therefore

\begin{equation*}
[p_ap_k][z^k]\bigl(z^{-a}E/Q\bigr)
=(-1)^{a+k-1}\bigl(t_k(A)+t_k(C)\bigr).
\tag{5.7}
\end{equation*}


The range $k<\min(a,b)$ is sufficient: when it ends because of $a$, all endpoint thresholds are
already zero; when it ends because of $b$, the remaining thresholds equal the known global ones.
Thus the merged endpoint arm multiset is recovered, and Lemma 5.1 recovers $\{A,C\}$. The middle
block is the complement in the global twig multiset, and Lemma 4.1 recovers the gaps. Hence the
equal-terminal-order case is complete.

\section{A unique minimum endpoint with known root}


Assume $a<c$ and that rooted $A$ is known. This is automatic when $A$ has at least three arms. It
also holds for a path endpoint outside the adjacent boundary discussed in Section 7.

If the left arrival event occurs first, (3.6) recovers $\alpha$. The second event involves $C$ and
$D=A\mathbin{-}_\alpha B$ from (3.8). When $c\neq d$, it first gives only the smaller component,
not the unordered pair. We record the label-separation argument because it prevents an otherwise
invalid use of the two-branch reconstruction theorem.

\Needspace{0.13\textheight}
\subsection*{Lemma 6.1 (one-sided second-event labels)}


Suppose two candidates have the same rooted $A$, the same gaps, global arm multiset, degree
sequence, and path sequence. A strict second-event exchange $C\cong D'$ with $c<d$ and $d'<c'$
is impossible.

\paragraph{Proof}


Let $U$ be the arm multiset of $A$. Since $D'$ is starlike, $B'$ has one arm, say of length $t$.
Spider reconstruction gives

\begin{equation*}
\Lambda_C=U\mathbin{\uplus}\{t+\alpha\}.
\tag{6.1}
\end{equation*}


The degree sequence forces $B$ also to have one arm, say of length $v$, and $C'$ to have
$|U|+1$ arms. Equality of the global arm multisets reduces the strict case to

\begin{equation*}
\begin{aligned}
\Lambda_A&=U_0\mathbin{\uplus}\{t\},
&\Lambda_B&=\{v\},
&\Lambda_C&=U_0\mathbin{\uplus}\{t,t+\alpha\},\\
\Lambda_{B'}&=\{t\},
&\Lambda_{C'}&=U_0\mathbin{\uplus}\{v,t+\alpha\},
\end{aligned}
\tag{6.2}
\end{equation*}


with $U_0\neq\varnothing$. Writing $H_0=\sum_{\ell\in U_0}h_\ell$, the path-sequence difference
factors as

\begin{equation*}
W_T-W_{T'}=(h_v-h_t)(1-x^\beta)(x^\alpha-1)H_0.
\tag{6.3}
\end{equation*}


All four factors are nonzero in the strict case, a contradiction. The reverse exchange is the same
after interchanging the candidates. $\square$

When $c\neq d$, Lemma 6.1 fixes the side of the recovered smaller component. Its order and the
known sum $b+c$ then fix the scalar layout. If the recovered component is $D$, the two-branch
result [8] reconstructs its unrooted type. The already rooted block $A$ locates the endpoint; an
alternative occurrence either gives the same decomposition or reverses the whole tree. If the
recovered component is $C$, a nonpath $C$ has its root at its unique branch vertex. A path $C$ is
rooted by Lemma 6.3 below.

The equal-order event $c=d$ needs a genuine sum-separation argument. For an arm multiset $U$, use
$e_i(U)$ as in Section 2. Let $z=x^\alpha$, $h=h_\alpha$, and let $D$ be the two-branch tree formed
from $A,B$. A direct classification of three-leaf subtrees gives

\begin{equation*}
J_D=e_3(A)+e_3(B)+h\bigl(e_2(A)+e_2(B)\bigr)
+z\bigl(e_2(A)e_1(B)+e_1(A)e_2(B)\bigr).
\tag{6.4}
\end{equation*}


\Needspace{0.13\textheight}
\subsection*{Lemma 6.2 (equal-order $P/Q$ exchange)}


Fix rooted $A$ and $\alpha$. If two allocations of the arms between $B$ and $C$ have the same
global arm multiset and the same path sequence at $c=d$, then either the allocations agree or there
are multisets $P,Q$ with equal total length such that

\begin{equation*}
\begin{aligned}
B&=M\mathbin{\uplus}P,&C&=S\mathbin{\uplus}Q,\\
B'&=M\mathbin{\uplus}Q,&C'&=S\mathbin{\uplus}P,
\end{aligned}
\tag{6.5}
\end{equation*}


where, if $A$ has arms $a_1,\ldots,a_r$,

\begin{equation*}
M=\{\underbrace{\alpha,\ldots,\alpha}_{r-1\text{ copies}}\},\qquad
S=\{\alpha+a_1,\ldots,\alpha+a_r\}.
\end{equation*}


Moreover equality of the three-leaf subtree data forces $P=Q$.

\paragraph{Proof}


Let $a,b,c,u$ denote the arm-sum polynomials of $A,B,C,B\uplus C$, and put

\begin{equation*}
H=x^\alpha f_A+h_{\alpha-1}.
\end{equation*}


After deleting terms fixed by $A,\alpha,u$, the path sequence of $D\cup C$ is

\begin{equation*}
K+b^2+(H-u)b.
\tag{6.6}
\end{equation*}


Equality for $b,b'$ gives

\begin{equation*}
(b-b')(b+b'+H-u)=0.
\tag{6.7}
\end{equation*}


The nontrivial factor, together with
$h_{\alpha+\ell}=h_\alpha+x^\alpha h_\ell$, gives the arm-multiplicity identity

\begin{equation*}
b'=c-S+M,
\tag{6.8a}
\end{equation*}


where the same letters denote the corresponding formal sums in the free abelian group on positive
arm lengths. Every element of $S$ is strictly larger than $\alpha$, whereas $M$ consists only of copies
of $\alpha$. Since $b'$ has nonnegative multiplicities, (6.8a) forces $C$ to contain every
occurrence of $S$. Moreover

\begin{equation*}
c'=u-b'=b+S-M.
\end{equation*}


Since $S$ has no support at length $\alpha$ and $c'$ has nonnegative multiplicities, this identity
forces $B$ to contain every occurrence of $M$. The nonnegative remainders are respectively $Q$ and
$P$, which gives exactly (6.5). The equality $c=d$ then yields $\sum P=\sum Q$.

For the three-leaf term, expand (6.4) and the spider term (2.6). All quadratic allocation terms
cancel, leaving

\begin{equation*}
(J_D+J_C)-(J_{D'}+J_{C'})
=x^\alpha(1-x^\alpha)e_2(A)\bigl(e_1(P)-e_1(Q)\bigr).
\tag{6.8}
\end{equation*}


The first three factors are nonzero. Equality therefore gives $e_1(P)=e_1(Q)$, and injectivity of
the arm-sum encoding gives $P=Q$. $\square$

The remaining possibility is that the right endpoint event $c$ occurs before the left chain reaches
the middle root, so $a<c<a+\alpha$. The derivative first gives $X_C$. If $C$ is a path, write
$\delta=c-a$ and subtract the known pruning term at the next order. The remainder is

\begin{equation*}
R=eX_{C^1}-fX_{A^{\delta+1}},
\qquad
e=\mathbf1_{\beta\geq2},\quad f=\mathbf1_{\alpha=\delta+1}.
\tag{6.9}
\end{equation*}


Principal specialization distinguishes the boundary states except possibly when
$X_{C^1}=X_{A^{\delta+1}}$. If $A$ is not a path, the two spiders have different degree sequences.
If $A$ is a path, this is the special collision treated in Section 7. Once the layout is fixed, the
following factorization roots $C$.

\Needspace{0.13\textheight}
\subsection*{Lemma 6.3 (fixed-layout path rooting)}


Suppose $A$ and the scalar layout are known, $C$ is an unrooted path, and $c<d$. Then the path
sequence determines the root position of $C$.

\paragraph{Proof}


Let $a(x),b(x),c(x)$ be the arm sums and set $u=b+c$, $z=x^\alpha$, $w=x^\beta$,
$h=h_\alpha$, $g=h_\beta$. After fixed terms are removed, the path-sequence remainder is

\begin{equation*}
\begin{aligned}
W={}&a(h+zg)+b(h+g)+c(g+wh)\\
&+(z-1)ab+(zw-1)ac+(w-1)bc.
\end{aligned}
\tag{6.10}
\end{equation*}


For a second allocation $b'=u-c'$, subtraction gives

\begin{equation*}
W-W'=(c-c')(w-1)(za+h+u-c-c').
\tag{6.11}
\end{equation*}


At $x=1$ the last factor is $d-c>0$. Thus equality forces $c=c'$, and the arm-sum injection fixes
the rooted path. $\square$

Lemmas 6.1--6.3, followed by Lemma 4.1, complete every uniquely minimal endpoint case for which the
root of $A$ is known.

\section{Rooting a uniquely minimal path endpoint}


Assume $a<c$ and $A$ is a two-arm path. If $\alpha\geq2$ and $c\geq a+2$, the coefficient of
order $a+1$ occurs before either event. If $\tau$ is the total number of twigs, then

\begin{equation*}
E_{a+1}=X_{A^1}+(\tau-2)p_1X_A.
\tag{7.1}
\end{equation*}


The unrooted $X_A$ and $\tau$ are known, so $X_{A^1}$ is extracted. The three-arm spider $A^1$ is
reconstructed, and deleting one added arm of length one roots $A$. The coefficient lies safely in
(3.4) because

\begin{equation*}
n-(a+1)=b+c+\alpha+\beta-3
\end{equation*}


exceeds every twig length. Only the adjacent boundary

\begin{equation*}
\alpha=1\quad\text{or}\quad c=a+1
\tag{7.2}
\end{equation*}


remains.

\Needspace{0.13\textheight}
\subsection*{7.1 Fixed scalar layout}


Fix $(a,b,c,\alpha,\beta)$. In (6.10), let the numbers of arms at the three centers be $2,q,m$.
For another allocation with the same global arm multiset, let $k\geq2$ be the first length threshold
at which the arm counts differ, and let the differences be $\rho,\sigma,\tau$. The global sequence,
the coefficient of $x^{k+1}$ in (6.10), and the coefficient of $x^{k+2}$ in the three-leaf
polynomial give

\begin{equation*}
\rho+\sigma+\tau=0,
\tag{7.3}
\end{equation*}


\begin{equation*}
-(m-2)\rho+(q-m+1)\sigma=0,
\tag{7.4}
\end{equation*}


\begin{equation*}
\left(1-\binom m2\right)\rho+
\left(\binom{q+1}2-\binom m2\right)\sigma=0.
\tag{7.5}
\end{equation*}


The determinant of (7.4)--(7.5) is

\begin{equation*}
-\frac{(m-2)(q-1)(q-m+1)}2.
\tag{7.6}
\end{equation*}


When it is nonzero, a first difference is impossible. In the three singular cases, the equations
recover the merged multiset of each equal-degree class. If $q=1$, the unique middle arm is fixed by
$b$; if $m=2$, the two endpoint paths have different total arm lengths $a-1<c-1$, which fixes the
pairing of their four arms; if $q=m-1$, the left endpoint lies in its own degree class. Hence rooted
$A$ is recovered.

The degree sequence also permits an apparent middle/right degree swap

\begin{equation*}
q'=m-1,\qquad m'=q+1.
\tag{7.7}
\end{equation*}


When $\alpha=1$, compare the coefficients $[x^3]W$, $[x^4]J$, and
$K_5=[x^5y^4]S_T$. If $A_2,A'_2$ count endpoint arms of length at least two, the last two
differences are

\begin{equation*}
\Delta J_4=
\frac{(A_2-A'_2)(q-1)(m-2)}2-(m-q-1),
\tag{7.8}
\end{equation*}


\begin{equation*}
\Delta K_5=
\frac{(A_2-A'_2)(m-2)(m+q)(q-1)}6
+\frac{(m+q)(q-m+1)}2.
\tag{7.9}
\end{equation*}


Here a five-edge, four-leaf subtree is either a subdivided four-star or a double three-star, so

\begin{equation*}
K_5=\sum_r\binom{\deg(r)-1}{3}|N_2(r)|
+\sum_{rs}\binom{\deg(r)-1}{2}\binom{\deg(s)-1}{2},
\tag{7.10}
\end{equation*}


where $N_2(r)$ is the set of vertices at distance two from $r$, and the second sum is over adjacent
branch vertices. If (7.8)--(7.9) vanish, elimination gives
$m=q+1$, so (7.7) is not an actual swap. When $c=a+1$ and $\alpha\geq2$, the derivative instead
extracts $X_{A^1}+X_C$; its degree sequence compares $\{3,m\}$ with $\{3,q+1\}$ and again forces
$m=q+1$. Thus a fixed scalar layout uniquely roots $A$.

\Needspace{0.13\textheight}
\subsection*{7.2 Classifying changes of scalar layout}


Delete one edge of $T$. By (2.1), the coefficient of $p_{n-s}p_s$ counts edges whose deletion has
component orders $\{s,n-s\}$. Subtracting the known twig-edge contributions leaves the trunk-cut
profile. In an oriented expansion it is the sum of four interval indicators

\begin{equation*}
\begin{aligned}
I_A&=[a,a+\alpha-1],&
I_C&=[c,c+\beta-1],\\
I_D&=[a+b+\alpha-1,a+b+\alpha+\beta-2],&
I_E&=[b+c+\beta-1,b+c+\alpha+\beta-2].
\end{aligned}
\tag{7.11}
\end{equation*}


Fix $n,a,L=\alpha+\beta$ and impose (7.2). Put $N=n+1$ and
$\phi(t)=x^t-x^{N-t}$. If $P(x)$ is the profile polynomial, then

\begin{equation*}
(1-x)P(x)=
\phi(a)-\phi(a+\alpha)+\phi(c)-\phi(c+\beta).
\tag{7.12}
\end{equation*}


This signed endpoint identity handles coincident endpoints and the half-order cut without separate
cases. If both candidates satisfy $c=a+1$, the variable part factors as

\begin{equation*}
(1-x^{b-1})(x^{a+\alpha}+x^{a+L+1-\alpha}),
\tag{7.13}
\end{equation*}


so equality permits only the original layout or the \textbf{gap exchange}

\begin{equation*}
(\alpha',\beta')=(\beta+1,\alpha-1).
\tag{7.14}
\end{equation*}


If both have $\alpha=1$, the variable part is

\begin{equation*}
(1-x^\beta)(x^c+x^{N-\beta-c}),
\tag{7.15}
\end{equation*}


so equality permits only the original layout or the \textbf{$C/D$ exchange}

\begin{equation*}
(b',c')=(c-a,a+b),\qquad\beta'=\beta.
\tag{7.16}
\end{equation*}


If the two candidates lie in different pure branches of (7.2), the coefficient of $x^{a+1}$ in
(7.12) has opposite signs. Thus (7.14)--(7.16) exhaust all cross-layout collisions.

\Needspace{0.13\textheight}
\subsection*{7.3 Closing the two exchanges}


For the $C/D$ exchange, the first-threshold equations (7.3)--(7.6) remain valid. If $q=1$ and
$m>2$, equations (7.4)--(7.5) force $\tau=0$. Hence $C=C'$ threshold by threshold, while the
three-arm multiset $A\mathbin{\uplus}B=A'\mathbin{\uplus}B'$ is fixed. Since $B,B'$ each consist
of one arm and $A,A'$ are two-arm paths of the same total length, two equal-sum two-submultisets of
the same three-element multiset cannot differ: if they share an occurrence, their other occurrences
agree, and disjointness is impossible. Thus $A=A'$.

If $m=2$ and $q>1$, equation (7.4) forces $\sigma=0$, so $B=B'$ and
$A\mathbin{\uplus}C=A'\mathbin{\uplus}C'$. Two different equal-sum pairs in the same four-element
multiset must be complements. That would make the arm sums of $A$ and $C$ equal, contradicting
$a<c$. Hence again $A=A'$. If $q=m-1>1$, equation (7.4) directly gives $\rho=0$, so threshold
induction gives $A=A'$. In their common intersection $q=1,m=2$, any nontrivial arm allocation has
the form

\begin{equation*}
\begin{aligned}
A&=\{p,t-p\},&B&=\{v\},&C&=\{t-v,\ell\},\\
A'&=\{v,t-v\},&B'&=\{p\},&C'&=\{t-p,\ell\},
\end{aligned}
\qquad \ell=p+v+1.
\tag{7.17}
\end{equation*}


With $P=x^p,V=x^v,T=x^t,w=x^\beta$, the path difference is

\begin{equation*}
W-W'=\frac{x^2(P-V)}{PV(x-1)^2}N,
\tag{7.18}
\end{equation*}


\begin{equation*}
N=(PV-T)w(1-x)(1-PVx)+T(1-P)(1-V)x(1-w).
\tag{7.19}
\end{equation*}


If $p\neq v$, the polynomial $N$ is nonzero: for $0<x<1$ its two terms are positive when
$t\geq p+v$; when $t<p+v$ its lowest term is positive, with $\beta=1$ handled after removing
$x(1-x)$. A degree-position swap is eliminated by the same pair (7.8)--(7.9). Once $A$ is fixed,
Lemma 6.1 rules out the remaining one-sided event-label exchange. Hence (7.16) is trivial.

For the gap exchange, put $r=\alpha-1<s=\beta$. At order $a+1=c$, (3.4) extracts

\begin{equation*}
X_{A^1}+X_C+(\tau-2)p_1X_A.
\tag{7.20}
\end{equation*}


The last term is known. For a $d$-arm spider, the highest-leaf subtree term is

\begin{equation*}
[y^d]S_T(x,y)=\prod_{\ell\in\Lambda}h_\ell(x),
\tag{7.21}
\end{equation*}


which recovers the arm multiset by cyclotomic divisibility and downward inversion. Except when the
right block has three arms, (7.20)--(7.21) already separate $A^1$ and $C$. Write
$\mathsf a,\mathsf b,\mathsf c$ for the arm-sum polynomials of $A,B,C$. The path difference then
factors as

\begin{equation*}
\Delta W=\mathsf b(x^r-x^s)\bigl(x(\mathsf a+1)-\mathsf c\bigr),
\tag{7.22}
\end{equation*}


whose last factor has $x$-coefficient $1-m\neq0$.

It remains to record why the three-arm degeneracies cannot cancel (7.22). Let $k$ be the first
threshold at which the assignments differ, with differences $\rho,\sigma,\tau$. Before the shorter
gap exponent $r$ is reached, the path and three-leaf equations are (7.4)--(7.5). If the assignments
remain equal until the scalar exchange appears, the first path contribution is $q(1-m)$. At the
border $k=r+1$, the next three-leaf contribution is

\begin{equation*}
-q\binom m2-\binom{q+1}2(m-1).
\tag{7.23}
\end{equation*}


For $m=3,q\geq3$, simultaneous cancellation would give
$\rho=2q/(q-1)>2$, impossible for a two-arm $A$. If $q=1$, the single middle arm is fixed and the
two components in (7.20) are either fixed or exchanged. In the exchanged case, write $U,V$ for
the arm-sum polynomials of the two candidate left path blocks. The exact path difference is

\begin{equation*}
(x-1)\mathsf b(x^rU-x^sV)+x(U-V)(x^{r+s+1}-1),
\tag{7.24}
\end{equation*}


whose first possible coefficient has absolute contribution two against at most one from the second
term. If $q=2$, let $D(x)$ be the difference between the arm-count polynomials of the two event
forests and put $N_0=a$. Expanding the three factors in (7.21) gives the reciprocity

\begin{equation*}
D(x)=x^{N_0}D(x^{-1}).
\end{equation*}


The global arm multiset says that the same $D$ is the negative of the difference between the two
middle two-arm blocks, hence also

\begin{equation*}
D(x)=x^{b-1}D(x^{-1}).
\end{equation*}


If $b-1\neq N_0$, the two identities force $D=0$ by comparing the least and greatest nonzero
degrees, reducing to the already separated fixed-component or exchanged-component cases. Thus the
only remaining possibility is

\begin{equation*}
m=3,\qquad b=c=a+1.
\tag{7.25}
\end{equation*}


Here the same seven arms are partitioned into a two-set of sum $a-1$, a two-set of sum $a$, and a
three-set of sum $a$. At the first changed threshold the two-set of sum $a-1$ must change; otherwise
complementary arm transfer forces an earlier change. Thus $\rho\neq0$. If $k\leq r$, the path
coefficient is $-\rho$; if $k\geq r+2$, the scalar term is $-4$; if $k=r+1$, their sum is
$-\rho-4$, nonzero since $|\rho|\leq2$. This closes (7.25).

Finally consider the apparent collision in (6.9) when both $A$ and $C$ are paths. Spider
reconstruction forces

\begin{equation*}
A=\{1,\eta\},\qquad C=\{\eta,\delta+1\},
\tag{7.26}
\end{equation*}


and the competing layouts are

\begin{equation*}
(\alpha,\beta)=(\delta+1,s),\qquad
(\alpha',\beta')=(s+\delta,1).
\tag{7.27}
\end{equation*}


Let the first candidate have $q$ middle arms; let the second candidate have $q'$ middle arms and
$m'$ right arms. Equality of the global arm count gives

\begin{equation*}
q+2=q'+m'.
\end{equation*}


Let $c_2,c'_2$ be the numbers of arms of length at least two in $C,C'$, and let $u_2$ be the
number of such arms in $B\mathbin{\uplus}C$. The coefficient of $x^2$ in the path difference is

\begin{equation*}
(2-m')(m'-q-1).
\tag{7.28}
\end{equation*}


When this vanishes, either $m'=2$ or $m'=q+1$. In the first case $q'=q$, and the coefficient of
$x^3$ is

\begin{equation*}
(q-1)(c'_2-c_2)-q.
\tag{7.29}
\end{equation*}


Both $C$ and $C'$ are two-arm paths of order at least five, so $c_2,c'_2\in\{1,2\}$. If
$c'_2-c_2\leq0$, (7.29) is negative; if the difference is one, it equals $-1$. Thus it is never
zero.

In the second case $q'=1$, and the coefficient of $x^3$ is

\begin{equation*}
(q-1)(u_2-c_2-c'_2)-1.
\tag{7.30}
\end{equation*}


The case $q=1$ already belongs to $m'=2$. If $q\geq3$, (7.30) is congruent to $-1$ modulo
$q-1$ and cannot vanish. It remains only to check $q=2$. Write the first candidate's middle arms as
$B=\{u,v\}$. The second candidate has the single middle arm $u+v$. Equality of the global arm
multisets says that this length occurs in

\begin{equation*}
\{u,v,\eta,\delta+1\}.
\end{equation*}


It is larger than $u$ and $v$, so it equals $\eta$ or $\delta+1$. If $u+v=\eta$, then
$C'=\{u,v,\delta+1\}$; if $u+v=\delta+1$, then $C'=\{u,v,\eta\}$. In either case
$u_2-c_2-c'_2\leq0$, so (7.30) is at most $-1$. Hence the right-first cross-layout collision is
also impossible.

Sections 7.1--7.3 prove that a uniquely smallest path endpoint is rooted and that its scalar layout
is unique. Lemmas 6.1--6.3 recover the remaining blocks, and Lemma 4.1 recovers the gaps.

\section{Proof of the main theorem}


Let $F$ be a finite simple graph with $X_F=X_T$. The two graphs have the same vertex count, edge
count, and girth. Since $T$ is a tree, $F$ is acyclic and has $|V(F)|-1$ edges; hence $F$ is
connected and is itself a tree. Their degree sequences agree, so $F$ also has exactly three branch
vertices. Their trunk orders and global twig multisets agree by [3].

Apply Lemma 3.1 to both trees. It recovers the minimum endpoint order and the sum of the endpoint
chromatic symmetric functions attaining it.

If the endpoint orders are equal, Section 5 recovers the unordered rooted pair of endpoint blocks
for every pair of gaps, including rooted-path cases and $\alpha=\beta=1$. The middle block is the
complement in the global twig multiset, and Lemma 4.1 recovers the two gaps up to global reversal.

If the minimum endpoint is unique, orient the display so that it is $A$. When $A$ has at least three
arms, its root is the unique branch vertex and Section 6 reconstructs $B,C,\alpha,\beta$. When $A$
is a path, equation (7.1) handles the nonadjacent range, while Section 7 handles exactly the
remaining boundary $\alpha=1$ or $c=a+1$ and the possible right-first collision. In every case the
three rooted blocks and both gaps are recovered, with global reversal as the only ambiguity.

Thus $F$ and $T$ have identical rooted block data and trunk placement, so $F\cong T$. $\square$

\section{Exact controls and proof boundary}


The proof is structural and applies to arbitrary positive arm and trunk lengths. Seventeen exact
Python controls accompany the proof-development record. They independently expand (2.1), verify
the derivative and event identities, enumerate small scalar layouts, check the threshold formulas,
and include destructive tests for two invalid shortcuts: extending a two-spider lemma verbatim to a
path--spider pair, and treating the first recovered second-event component as if the complete
unordered component pair were already known.

The final regression run passed all seventeen controls. The largest classification guard examined
$1,215,450$ scalar descriptions, and the one-sided label-swap guard checked $16,320$ factor cases.
These computations test formulas and boundary coverage only. The conclusion of Theorem 1.1 rests on
the general arguments in Sections 3--8, not on finite enumeration.

The theorem closes the natural layer of exactly three branch vertices. It does not cover arbitrary
trees or trees with four branch vertices. With four branch vertices the branch core need not be a
path, and the three-block allocation and two-gap exchange classification used here no longer
exhaust the possibilities.

\section{Relation to prior work}


The closest direct predecessor is Wang, Yu and Zhang [8], who proved the exact-two-branch layer and
suggested that their path-trunk method might extend to three. The present proof uses their weighted
deletion--contraction and two-spider reconstruction results, but the full three-branch layer requires
new endpoint-event, path-rooting, threshold-moment, and exchange-separation arguments.

Other known families overlap the theorem without covering it. Caterpillars [4], 2-spiders [5],
proper $q$-caterpillars [6], and diameter-at-most-five trees [7] each include some trees with three
branch vertices. The multiplicity-isolated proper-tree criteria of Zeng and Awan [10,11] cover
further subfamilies under weight-separation hypotheses. Theorem 1.1 imposes none of those subclass
conditions, but conversely does not contain all trees treated by those works. Aliste-Prieto, Martin,
Wagner and Zamora [9] provide the generalized-degree and subtree-polynomial framework used to
identify which weaker polynomial data can safely be extracted from $X_T$. Loehr and Warrington's
rooted chromatic symmetric function distinguishes rooted trees [12], but it is an additional rooted
invariant; the rooting steps in Sections 5--7 are instead extracted from the ordinary $X_T$.

\section*{References}
\begin{enumerate}
\footnotesize

\item R. P. Stanley, ``A symmetric function generalization of the chromatic polynomial of a graph,'' \emph{Advances in Mathematics} \textbf{111} (1995), 166--194. \url{https://doi.org/10.1006/aima.1995.1020}
\item J. L. Martin, M. Morin, and J. D. Wagner, ``On distinguishing trees by their chromatic symmetric functions,'' \emph{Journal of Combinatorial Theory, Series A} \textbf{115} (2008), 237--253. \url{https://doi.org/10.1016/j.jcta.2007.05.008}
\item L. Crew, ``A note on distinguishing trees with the chromatic symmetric function,'' \emph{Discrete Mathematics} \textbf{345} (2022), 112682. \url{https://doi.org/10.1016/j.disc.2021.112682}
\item M. Loebl and J.-S. Sereni, ``Isomorphism of weighted trees and Stanley's isomorphism conjecture for caterpillars,'' \emph{Annales de l'Institut Henri Poincaré D} \textbf{6} (2019), 357--384. \url{https://doi.org/10.4171/aihpd/74}
\item J. Huryn and S. Chmutov, ``A few more trees the chromatic symmetric function can distinguish,'' \emph{Involve} \textbf{13} (2020), 109--116. \url{https://doi.org/10.2140/involve.2020.13.109}
\item A. Ganesan, N. Narayanan, B. V. R. Rao, and S. S. Sawant, ``Proper $q$-caterpillars are distinguished by their chromatic symmetric functions,'' \emph{Discrete Mathematics} \textbf{347} (2024), 114162. \url{https://doi.org/10.1016/j.disc.2024.114162}
\item M. Gonzalez, R. Orellana, and M. Tomba, ``The chromatic symmetric function in the star-basis,'' \emph{Discrete Mathematics} \textbf{349} (2026), 114691. \url{https://doi.org/10.1016/j.disc.2025.114691}
\item Y. Wang, X. Yu, and X.-D. Zhang, ``A class of trees determined by their chromatic symmetric functions,'' \emph{Discrete Mathematics} \textbf{347} (2024), 114096. \url{https://doi.org/10.1016/j.disc.2024.114096}
\item J. Aliste-Prieto, J. L. Martin, J. D. Wagner, and J. Zamora, ``Chromatic symmetric functions and polynomial invariants of trees,'' \emph{Bulletin of the London Mathematical Society} \textbf{56} (2024), 3452--3476. \url{https://doi.org/10.1112/blms.13144}
\item Z. Zeng, ``Multiplicity-isolated cores and chromatic symmetric reconstruction of trees,'' arXiv preprint, 2026. \url{https://arxiv.org/abs/2608.18179}
\item S. A. Awan, ``First-derivative chromatic symmetric reconstruction for proper trees,'' arXiv preprint, 2026. \url{https://arxiv.org/abs/2609.13464}
\item N. A. Loehr and G. S. Warrington, ``A rooted variant of Stanley's chromatic symmetric function,'' \emph{Discrete Mathematics} \textbf{347} (2024), 113805. \url{https://doi.org/10.1016/j.disc.2023.113805}

\end{enumerate}

\section*{AI-assisted research and writing disclosure}


Carptopus is the author and is responsible for the mathematical claims, source selection, and
public release. OpenAI Codex was used as an AI-assisted tool for symbolic exploration, exact
verification, adversarial proof review, manuscript drafting, and editing.

\section*{License}


This manuscript is licensed under the \href{https://creativecommons.org/licenses/by/4.0/}{Creative Commons Attribution 4.0 International License (CC BY 4.0)}.

\end{document}
